\documentclass[../../../thesis.tex]{subfiles}
\begin{document}
\subsection{Záverečná časť}
Na záver práce uvádzame dodatky a prílohy.
Prílohy práce sú zväčša materiály,
ktoré majú odlišný formát voči samotnej práci.
Sú to napríklad pamäťové nosiče,
dátové súbory, veľkoformátové mapy, výkresy a podobne.
Každú prílohu treba jasne označiť, očíslovať a~nazvať.
Zoznam príloh potom uvedieme v~jednom z~dodatkov.

Do tzv. dodatkov umiestňujeme informácie,
ktoré kvôli rozsahu nemôžu byť v hlavnom texte práce.
Sú to napríklad údajové listy k~použitým prístrojom
a~zariadeniam, zdĺhavejšie matematické odvodenia,
rozsiahlejšie kódy programov, dokumentácia
k~vytvoreným programom, definície neštandardných objektov,
ktoré v práci používame,
série rozsiahlych výsledkov alebo meraní
a~ich grafy, fotografie a podobne.

Jednotlivé kapitoly v~dodatkoch číslujeme veľkými písmenami, čísla podkapitol majú formu A.1, B.3.2, atď.
Na tento účel vytvoríme pre každý dodatok samostatný súbor v~priečinku \verb|includes/|, odporúčame názov súboru v~tvare \verb|attachmentA.tex| alebo podobne.
Každý dodatok je potom potrebné načítať v~hlavnom súbore \verb|thesis.tex| nasledujúcim spôsobom:
\begin{verbatim}
\FEIappendix{Názov prílohy\label{att:A}}{includes/attachmentA}
\end{verbatim}
Prvý parameter makra je názov dodatku a ten sa nesmie nachádzať v~zdrojovom súbore \verb|attachemntA.tex|.
\end{document}